\chapter{Výskumná štúdia}
\label{chap:vyskum}

Návrh hodnotiaceho systému vychádza z predpokladov o ľudskom vnímaní prezentácií: že jednotlivé dimenzie majú rôzny vplyv na celkový dojem a že sústredené výskyty negatívnych javov sú pre publikum relevantnejšie ako rovnomerne roztrúsené.
Výskumná štúdia tieto predpoklady testuje --- jej výsledky spätne validujú konkrétne dizajnové rozhodnutia systému a umožňujú kalibrovať hodnotiaci model na základe skutočného ľudského vnímania.

Prepojenie medzi hypotézami a systémom zachytáva nasledujúca tabuľka:

\begin{table}[H]
    \centering
    \begin{tabularx}{\textwidth}{|l|X|}
        \hline
        \textbf{Hypotéza} & \textbf{Čo validuje v systéme} \\
        \hline
        \specialrule{1.2pt}{0pt}{0pt}
        H1.1a & Oprávnenosť toho, že vizuálne a hlasové dimenzie majú dohromady vyššiu váhu ako plynulosť reči \\ \hline
        H1.1b & Oprávnenosť hierarchie váh: vizuálne $>$ hlasové $>$ plynulosť \\ \hline
        H1.2  & Oprávnenosť vyššej váhy hlasu oproti plynulosti reči \\ \hline
        H2.1  & Oprávnenosť segmentačného prístupu. V tejto štúdii sa zameriavam iba na výplňové slová \\ \hline
    \end{tabularx}
    \caption{Prepojenie výskumných hypotéz s dizajnovými rozhodnutiami systému}
    \label{tab:hypothesesSystem}
\end{table}

\section{Výskumné otázky a hypotézy}

Štúdia sa zameriava na dve výskumné otázky:

\textbf{RQ1:} Majú vizuálne dimenzie (očný kontakt, pohyb tela) a hlasové dimenzie (výška hlasu, hlasitosť) väčší vplyv na vnímané celkové hodnotenie prezentácie ako plynulosť reči (výplňové slová)? Platí hierarchia negatívneho efektu na dojem z prezentácie týchto dimenzií: vizuálne~>~hlasové~>~plynulosť?

\textbf{RQ2:} Je sústredený výskyt negatívneho javu v jednom súvislom bloku vnímaný respondentmi horšie ako rovnaký celkový čas javu rozloženého rovnomerne počas celej prezentácie?

Na základe týchto otázok boli definované štyri hypotézy:

\begin{itemize}
    \item \textbf{H1.1a} --- Očný kontakt, pohyby tela a vlastnosti hlasu majú spolu väčší vplyv na celkové hodnotenie prezentácie ako plynulosť reči.
    \item \textbf{H1.1b} --- Očný kontakt a pohyby tela majú väčší vplyv na celkové hodnotenie prezentácie ako hlasové vlastnosti rečníka, ktoré sú zároveň vplyvnejšie ako výplňové slová.
    \item \textbf{H1.2} --- Monotónnosť hlasu má väčší negatívny vplyv na celkové hodnotenie prezentácie ako výplňové slová.
    \item \textbf{H2.1} --- Sústredený výskyt výplňových slov má väčší negatívny vplyv na celkové hodnotenie prezentácie ako ich rovnomerne roztrúsený výskyt.
\end{itemize}

\section{Dizajn experimentu}

\subsection{Príprava videí}

Pre štúdiu bola nahratá séria šiestich videí s rovnakým prezentujúcim a rovnakým obsahom. Každé video malo dĺžku približne 60 sekúnd. Videá sa líšili práve tým aspektom prezentácie, ktorý bol predmetom skúmania. Manipulácia vždy zasiahla iba jednu dimenziu, ostatné boli ponechané neutrálne.

\begin{table}[H]
    \centering
    \begin{tabularx}{\textwidth}{|l|l|X|}
        \hline
        \textbf{Video} & \textbf{Označenie} & \textbf{Popis manipulácie} \\
        \hline
        \specialrule{1.2pt}{0pt}{0pt}
        V1 & Baseline & Neutrálna prezentácia bez zámerných chýb. Slúži ako referenčná hodnota. \\
        \hline
        V2 & Zlý očný kontakt & Prezentujúci sa pozerá na poznámky a mimo kamery --- simulácia absencie očného kontaktu. \\
        \hline
        V3 & Pohyb bokov & Výrazné kývanie bokov zo strany na stranu počas celej prezentácie. \\
        \hline
        V4 & Monotónny hlas & Prezentujúci hovorí s konštantnou výškou a rytmom hlasu, bez akejkoľvek prozódie. \\
        \hline
        V5 & Výplňové slová (rovnomerne) & Výplňové slová sú rozložené rovnomerne počas celej prezentácie. \\
        \hline
        V6 & Výplňové slová (sústredene) & Rovnaký celkový počet výplňových slov sústredený do strednej časti prezentácie. \\
        \hline
    \end{tabularx}
    \caption{Prehľad videí použitých v štúdii}
    \label{tab:videos}
\end{table}

\subsection{Zber dát}

Dáta boli zozbierané prostredníctvom online dotazníka (Google Forms).
Respondentom bolo postupne zobrazených všetkých šesť videí v rovnakom poradí.
Po každom videu hodnotili prezentujúceho na škále 0--10 v piatich kategóriach: \textbf{celkové hodnotenie}, \textbf{očný kontakt}, \textbf{pohyb tela}, \textbf{plynulosť reči} a \textbf{hlas}.

Celkový počet respondentov bol \textbf{n~=~52}.

Demografické charakteristiky vzorky sú uvedené v tabuľkách~\ref{tab:demographics} a~\ref{tab:experience}. Vzorka zahŕňala 16 žien (30,8\,\%) a 36 mužov (69,2\,\%). Väčšinu tvorili respondenti vo vekovej kategórii 18--25 rokov (46,2\,\%). Z hľadiska skúseností s prezentovaním dominovali respondenti, ktorí prezentujú len výnimočne (44,2\,\%), nasledovaní tými, pre ktorých je prezentovanie bežnou súčasťou štúdia alebo práce (34,6\,\%). Vzorka pokrývala rôzne úrovne skúseností, čo zvyšuje jej reprezentatívnosť.

\begin{table}[H]
    \centering
    \begin{tabular}{|l|c|c|}
        \hline
        \textbf{Charakteristika} & \textbf{Počet} & \textbf{Podiel} \\
        \hline
        \specialrule{1.2pt}{0pt}{0pt}
        \multicolumn{3}{|l|}{\textit{Pohlavie}} \\ \hline
        Ženy & 16 & 30,8\,\% \\ \hline
        Muži & 36 & 69,2\,\% \\ \hline
        \multicolumn{3}{|l|}{\textit{Veková skupina}} \\ \hline
        18--25 rokov & 24 & 46,2\,\% \\ \hline
        26--35 rokov & 12 & 23,1\,\% \\ \hline
        36--50 rokov & 14 & 26,9\,\% \\ \hline
        51 a viac rokov & 2 & 3,8\,\% \\ \hline
    \end{tabular}
    \caption{Demografické charakteristiky respondentov}
    \label{tab:demographics}
\end{table}

\begin{table}[H]
    \centering
    \begin{tabular}{|l|c|c|}
        \hline
        \textbf{Skúsenosti s prezentovaním} & \textbf{Počet} & \textbf{Podiel} \\
        \hline
        \specialrule{1.2pt}{0pt}{0pt}
        Prezentovanie je bežná súčasť štúdia/práce & 18 & 34,6\,\% \\ \hline
        Prezentuje výnimočne & 23 & 44,2\,\% \\ \hline
        Dlhšie neprezentoval & 10 & 19,2\,\% \\ \hline
        Aktívne sa venuje prezentovaniu & 1 & 1,9\,\% \\ \hline
    \end{tabular}
    \caption{Skúsenosti respondentov s prezentovaním}
    \label{tab:experience}
\end{table}

\section{Štatistické metódy}

Pre overenie hypotéz boli použité nasledujúce štatistické metódy, implementované v jazyku Python s využitím knižnice \texttt{pingouin}\cite{pingouin}:

\begin{itemize}
    \item \textbf{Viacnásobná lineárna regresia} --- štandardizované beta koeficienty vyjadrujúce relatívny vplyv každej dimenzie na celkové hodnotenie. Dáta boli pred regresiou štandardizované (z-score).
    \item \textbf{Pearsonova korelácia per respondent} --- pre každého respondenta bola vypočítaná korelácia medzi hodnoteniami dimenzií a celkovým hodnotením naprieč šiestimi videami. Takto vzniklo 52 párov hodnôt pre každú dimenziu.
    \item \textbf{Párový t-test} --- porovnanie priemerov dvoch skupín merané na tých istých respondentoch. Použitý pre H1.1, H1.2 a H2.1.
    \item \textbf{Wilcoxonov test} --- neparametrická alternatíva k párovému t-testu. Slúži ako robustnostná kontrola pri porušení normality.
    \item \textbf{Cohenovo d} --- veľkosť efektu (malý: $|d| < 0{,}2$, stredný: $|d| \approx 0{,}5$, veľký: $|d| > 0{,}8$).
    \item \textbf{Repeated Measures ANOVA} s Greenhouse-Geisserovou korekciou --- testuje, či existujú štatisticky významné rozdiely medzi všetkými šiestimi videami súčasne. Post-hoc párové testy s Bonferroniho korekciou identifikujú konkrétne rozdielne dvojice.
    \item \textbf{Spearmanova rank korelácia} --- porovnanie poradia videí podľa modelu s poradím podľa ľudských hodnotení.
\end{itemize}

Pred spustením RM ANOVA boli overené jej predpoklady: normalita pomocou Shapiro-Wilkovho testu, sféricita pomocou Mauchlyho testu a prítomnosť odľahlých hodnôt podľa IQR pravidla. Normalita bola mierne porušená pre V1 ($p < 0{,}001$) a V5 ($p = 0{,}047$), čo je pri $n = 52$ akceptovateľné vzhľadom na Centrálnu limitnú vetu. Sféricita bola porušená ($W = 0{,}476$, $p < 0{,}001$, $\varepsilon = 0{,}780$), preto bola použitá Greenhouse-Geisserova korekcia.

\section{Výsledky}

Tabuľka~\ref{tab:hypotheses_summary} poskytuje stručný prehľad všetkých testovaných hypotéz a ich výsledkov. Podrobné štatistiky a interpretácia sú uvedené v nasledujúcich podkapitolách.

\begin{table}[H]
    \centering
    \begin{tabularx}{\textwidth}{|l|X|l|l|}
        \hline
        \textbf{Hypotéza} & \textbf{Čo testuje} & \textbf{Výsledok} & \textbf{p-hodnota} \\
        \hline
        \specialrule{1.2pt}{0pt}{0pt}
        H1.1a & Vizuálne+hlasové majú väčší vplyv ako plynulosť & Zamietnutá & 0,037 \\ \hline
        H1.1b & Hierarchia: Vizuálne > Hlasové > Plynulosť & Čiastočná & $<$0,001 \\ \hline
        H1.2 & Monotónnosť vnímaná horšie ako fillery & Nepodporená & 0,870 \\ \hline
        H2.1 & Sústredené fillery horšie ako rovnomerne & Zamietnutá & $<$0,001 \\ \hline
    \end{tabularx}
    \caption{Prehľad hypotéz a ich výsledkov}
    \label{tab:hypotheses_summary}
\end{table}

\subsection{H1.1 --- Vplyv dimenzií na celkové hodnotenie}

Ako prvý krok boli odhadnuté relatívne váhy jednotlivých dimenzií pomocou viacnásobnej lineárnej regresie. Dáta boli pred regresiou štandardizované (z-skóre), takže výsledné beta koeficienty sú priamo porovnateľné naprieč dimenziami:

\begin{table}[H]
    \centering
    \begin{tabular}{|l|c|}
        \hline
        \textbf{Dimenzia} & \textbf{Beta koeficient} \\
        \hline
        \specialrule{1.2pt}{0pt}{0pt}
        Hlas & 0,843 \\
        \hline
        Očný kontakt & 0,586 \\
        \hline
        Plynulosť & 0,574 \\
        \hline
        Pohyb tela & 0,250 \\
        \hline
    \end{tabular}
    \caption{Štandardizované beta koeficienty z lineárnej regresie}
    \label{tab:betas}
\end{table}

Empirická hierarchia vplyvu na celkové hodnotenie je teda \textbf{hlas~>~očný kontakt~>~plynulosť~>~pohyb tela}, čo je v rozpore s predpokladanou hierarchiou vizuálne~>~hlasové~>~plynulosť.

\textbf{H1.1a}

\textit{Nulová hypotéza:} Vizuálne a hlasové dimenzie spolu nie sú silnejším prediktorom celkového hodnotenia ako plynulosť reči --- teda rozdiely medzi skupinami sú náhodné.

\textit{Metóda:} Pre každého respondenta bola vypočítaná korelácia jeho hodnotení každej dimenzie s celkovým hodnotením naprieč 6 videami. Korelácie boli Fisherovou z-transformáciou prevedené do symetrického rozdelenia, čím sa umožnil párový t-test medzi skupinou (vizuálne + hlasové, priemerované) a plynulosťou. Tento prístup zohľadňuje individuálne rozdiely medzi respondentmi.

\textit{Výsledok:} Párový t-test: $t = -2{,}143$, $p = 0{,}037$. Priemerná korelácia skupiny vizuálne + hlasové ($\bar{r} = 0{,}550$) je štatisticky signifikantne \textbf{nižšia} ako korelácia plynulosti ($\bar{r} = 0{,}613$). \textbf{Hypotéza H1.1a bola zamietnutá v opačnom smere} --- vizuálne a hlasové dimenzie spolu nie sú silnejším prediktorom ako plynulosť.
Príčinou je nízky príspevok pohybu tela ($\beta = 0{,}250$), ktorý strháva priemer skupiny pod úroveň plynulosti.

\textbf{H1.1b}

\textit{Nulová hypotéza:} Medzi troma skupinami dimenzií (vizuálne, hlasové, plynulosť) nie je štatisticky významný rozdiel vo vplyve na celkové hodnotenie.

\textit{Metóda:} Rovnaký prístup ako pri H1.1a --- párový t-test na per-respondent koreláciách, tentokrát s dvoma porovnaniami: vizuálne vs. hlasové a hlasové vs. plynulosť.

\textit{Výsledok:} Vizuálne vs. hlasové: $t = -5{,}238$, $p < 0{,}001$ --- hlas štatisticky signifikantne dominuje nad vizuálnymi dimenziami. Hlasové vs. plynulosť: $t = 1{,}457$, $p = 0{,}151$ --- rozdiel nie je signifikantný. \textbf{Hypotéza H1.1b bola čiastočne podporená v opačnom smere} --- hlas dominuje nad vizuálnymi dimenziami (potvrdené), ale poradie hlas~>~plynulosť nebolo štatisticky dokázané.

\textit{Metodologická poznámka --- halo efekt:} Respondenti hodnotili každé video vo všetkých dimenziách súčasne. V dátach sa prejavil halo efekt: keď bolo video zhoršené v jednej dimenzii, klesali aj hodnotenia ostatných dimenzií, ktoré neboli manipulované (napr. V4 s monotónnym hlasom mal nižšie hodnotenie plynulosti 5,67 oproti baseline 8,18). Beta koeficienty teda čiastočne odrážajú aj tento efekt, nie len izolovaný vplyv každej dimenzie.

\subsection{H1.2 --- Monotónnosť vs. výplňové slová}

\textit{Nulová hypotéza:} Priemerné celkové hodnotenie videa s monotónnym hlasom (V4) sa štatisticky nelíši od hodnotenia videa s výplňovými slovami (V5).

\textit{Metóda:} Priame porovnanie celkových hodnotení V4 a V5 párovým t-testom --- keďže každý respondent hodnotil obe videá, eliminujú sa medziosobné rozdiely v prísnosti hodnotenia. Wilcoxonov test slúži ako robustnostná kontrola pri prípadnom porušení normality.

\textit{Výsledok:} Priemerné hodnotenie V4 = 4,40~±~1,91, V5 = 4,44~±~2,09. Párový t-test: $t = -0{,}164$, $p = 0{,}870$, $d = -0{,}019$. Wilcoxon: $W = 386{,}0$, $p = 0{,}742$. \textbf{Hypotéza H1.2 nebola podporená} --- oba javy boli respondentmi hodnotené prakticky identicky. Výplňové slová boli v štúdii prítomné vo vysokej frekvencii, čo pravdepodobne spôsobilo rovnako silný negatívny dojem ako monotónny hlas.

\subsection{H2.1 --- Distribúcia výplňových slov}

\textit{Nulová hypotéza:} Distribúcia výplňových slov v čase nemá vplyv na celkové hodnotenie --- prezentácie V5 (rovnomerne) a V6 (sústredene) budú hodnotené rovnako.

\textit{Metóda:} Párové porovnanie celkových hodnotení V5 a V6 párovým t-testom a Wilcoxonovým testom. Obe videá mali identický celkový počet výplňových slov, líšili sa iba ich rozmiestnením --- V5 rovnomerne po celej prezentácii, V6 sústredene v strednej časti.

\textit{Výsledok:} Priemerné hodnotenie V5 = 4,44~±~2,09, V6 = 6,27~±~2,07. Párový t-test: $t = 5{,}492$, $p < 0{,}001$, $d = 0{,}871$. Wilcoxon: $W = 81{,}0$, $p < 0{,}001$. 40 z 52 respondentov (76,9\,\%) hodnotilo V5 horšie ako V6. \textbf{Hypotéza H2.1 bola zamietnutá s opačným nálezom} --- rovnomerne rozložené výplňové slová sú vnímané výrazne horšie ako sústredené. Konštantná prítomnosť výplňových slov počas celej prezentácie neposkytuje poslucháčovi žiadnu pauzu od rušivého javu. Efekt je veľký ($d = 0{,}871$) a robustný naprieč oboma testmi.

\subsection{Repeated Measures ANOVA}

RM ANOVA potvrdila, že medzi videami existujú štatisticky vysoko významné rozdiely: $F(5, 255) = 39{,}412$, $p_{GG} < 10^{-22}$, $\eta^2_g = 0{,}256$ (veľký efekt). Bonferroniho post-hoc testy odhalili 12 z 15 signifikantných párov ($p < 0{,}05$). Nesignifikantné páry (V4 vs. V5, V2 vs. V6, V3 vs. V6) sú v súlade s ostatnými nálezmi.

RM ANOVA per dimenzia odhalila najväčší efekt pre \textbf{hlas} ($\eta^2_g = 0{,}402$), nasledovaný plynulosťou ($\eta^2_g = 0{,}334$) a očným kontaktom ($\eta^2_g = 0{,}175$). Pohyb tela mal najmenší efekt ($\eta^2_g = 0{,}050$). Tieto hodnoty sú v súlade s beta koeficientmi z regresie.

\section{Implikácie pre systém}

Výsledky štúdie boli zapracované do systému v dvoch krokoch:

\textbf{Model v2 --- aktualizácia váh:} Pôvodné váhy dimenzií (hlas 25\,\%, plynulosť 20\,\%, pohyb tela 25\,\%, očný kontakt 30\,\%) boli nahradené váhami odvodenými z normalizovaných beta koeficientov (súčet beta = 2,24):

\begin{table}[H]
    \centering
    \begin{tabular}{|l|c|c|}
        \hline
        \textbf{Dimenzia} & \textbf{Model v1} & \textbf{Model v2} \\
        \hline
        \specialrule{1.2pt}{0pt}{0pt}
        Hlas & 25\,\% & 37\,\% \\
        \hline
        Plynulosť & 20\,\% & 27\,\% \\
        \hline
        Očný kontakt & 30\,\% & 27\,\% \\
        \hline
        Pohyb tela & 25\,\% & 9\,\% \\
        \hline
    \end{tabular}
    \caption{Porovnanie váh dimenzií medzi modelmi}
    \label{tab:weights}
\end{table}

\textbf{Model v3 --- penalizácia distribúcie výplňových slov:} Na základe nálezu H2.1 bol do výpočtu skóre plynulosti pridaný distribučný trest. Prezentácia je rozdelená do 6 časových okien; koeficient variácie (CV = $\sigma / \mu$) hustoty výplňových slov určuje mieru rovnomernosti distribúcie. Nízke CV (rovnomerné rozloženie) vedie k vyššej penalizácii (max. 7 bodov).

\subsection{Porovnanie modelov}

Validita modelov bola overená Spearmanovou rank koreláciou medzi poradím 6 videí podľa systémového skóre a poradím podľa priemerných ľudských hodnotení. Spearmanova korelácia bola zvolená, pretože systémové skóre a ľudské hodnotenia sú na rôznych škálach (0--100 vs. 0--10), a preto nie je zmysluplné porovnávať absolútne hodnoty --- dôležité je, či systém zoraďuje videá rovnako ako ľudia. Rank korelácia je na tento účel priamo určená.

\begin{table}[H]
    \centering
    \begin{tabular}{|l|c|c|}
        \hline
        \textbf{Model} & \textbf{Spearman $\rho$} & \textbf{Zmena} \\
        \hline
        \specialrule{1.2pt}{0pt}{0pt}
        v1 (pôvodné váhy) & 0,600 & --- \\
        \hline
        v2 (váhy z výskumu) & 0,829 & $+0{,}229$ \\
        \hline
        v3 (váhy + distribúcia) & 0,829 & $0{,}000$ \\
        \hline
    \end{tabular}
    \caption{Spearmanová rank korelácia modelov s ľudskými hodnoteniami}
    \label{tab:spearman}
\end{table}

Model v2 dosiahol výrazné zlepšenie oproti modelu v1 --- poradie videí zodpovedá ľudskému vnímaniu výrazne presnejšie. Model v3 nezmenil celkové poradie (V5 bolo pred V6 už v modeli v2), avšak zväčšil rozdiel skóre V5 a V6 v dimenzii plynulosť z 21,8 na 28,8 bodu. Porovnaním tohto systémového rozdielu ($\Delta_\text{sys} = -2{,}88$) s reálnym rozdielom v ľudských hodnoteniach plynulosti ($\Delta_\text{human} = -1{,}64$, párový t-test: $t = -5{,}864$, $p < 0{,}001$) vyplýva, že distribučná penalizácia je mierne agresívnejšia, ako zodpovedá ľudskému vnímaniu --- systém trestá rovnomerne rozložené výplňové slová výraznejšie, ako to ľudia skutočne pociťujú. Kalibrácia tejto penalizácie ostáva predmetom ďalšieho vývoja.

Systémové skóre naprieč všetkými modelmi systematicky nadhodnocuje kvalitu prezentácie v porovnaní s ľudskými hodnoteniami (priemerný rozdiel $\approx +24$ bodov). Príčinou je, že systém deteguje len merateľné technické javy, zatiaľ čo ľudia vnímajú prezentáciu širšie --- vrátane aspektov, ktoré systém neanalyzuje (nervozita, dikcia, obsah). Kalibrácia absolútnych hodnôt skóre zostáva otvorenou výzvou pre ďalší vývoj.

\section{Validácia hodnotiaceho modelu}

Tretia výskumná oblasť overuje, či poradie videí podľa systémového skóre zodpovedá poradiu podľa ľudského hodnotenia a či empiricky odvodené váhy a distribučná penalizácia túto zhodu zlepšujú.

\subsection{H3.1 --- Výskumom odvodené váhy zvyšujú zhodu modelu s ľudským hodnotením}

Hypotéza: Spearmanova korelácia poradia videí bude vyššia pre model v2 (empirické váhy) ako pre model v1 (intuitívne váhy).

Výsledky sú zhrnuté v tabuľke~\ref{tab:spearman} v predchádzajúcej sekcii. Prechod z modelu v1 ($\rho = 0{,}600$, $p = 0{,}208$) na model v2 ($\rho = 0{,}829$, $p = 0{,}042$) predstavuje nárast o $+0{,}229$. Model v1 nie je štatisticky signifikantný, model v2 áno. \textbf{Záver: Hypotéza H3.1 potvrdená.}

Ako doplnková kontrola bola vykonaná leave-one-out krížová validácia (LOO-CV): pre každé z 6 videí boli váhy odvodené z regresie na zostávajúcich 5 videách a aplikované na vynechané video. Výsledná korelácia dosiahla $\rho = 0{,}829$ ($p = 0{,}042$, $\text{MAE} = 0{,}43$), čo naznačuje, že regresne odvodené váhy generalizujú aj mimo trénovacej vzorky.

\subsection{H3.2 --- Distribučná penalizácia presnejšie kopíruje ľudský rozdiel V5 vs. V6}

Hypotéza: Absolútna chyba $|\Delta_\text{sys} - \Delta_\text{human}|$ na dimenzii plynulosť bude nižšia pre model v3 ako pre modely v1 a v2.

Ľudský rozdiel hodnotenia plynulosti V5~$-$~V6: $\Delta_\text{human} = -1{,}64$ (párový t-test: $t = -5{,}864$, $p < 0{,}001$).

\begin{table}[H]
    \centering
    \begin{tabular}{|l|c|c|}
        \hline
        \textbf{Model} & \textbf{Systémový $\Delta$ (V5\,$-$\,V6)} & \textbf{$|\Delta_\text{sys} - \Delta_\text{human}|$} \\
        \hline
        \specialrule{1.2pt}{0pt}{0pt}
        v1 & $-2{,}18$ & 0,545 \\
        \hline
        v2 & $-2{,}18$ & 0,545 \\
        \hline
        v3 & $-2{,}88$ & 1,245 \\
        \hline
    \end{tabular}
    \caption{Porovnanie systémového a ľudského rozdielu V5~$-$~V6 na dimenzii plynulosť}
    \label{tab:h32}
\end{table}

\textbf{Záver: Hypotéza H3.2 nebola potvrdená --- výsledok je opačný.} Modely v1 a v2 sú výrazne bližšie k ľudskému rozdielu ($|\text{chyba}| = 0{,}545$) ako model v3 ($|\text{chyba}| = 1{,}245$). Distribučná penalizácia zveličuje rozdiel medzi rovnomerne a sústrednene rozloženými výplňovými slovami nad mieru, akú ľudia skutočne vnímajú. Smer efektu (V5 vnímané horšie ako V6) je zachytený správne vo všetkých troch modeloch, avšak veľkosť penalizácie v modeli v3 si vyžaduje ďalšiu kalibráciu.

\subsection{Obmedzenia štúdie}

\textbf{Fixné poradie videí.} Všetci respondenti videli videá v rovnakom poradí. To môže vytvárať order effects --- napríklad únava pri posledných videách alebo kontrastné efekty voči predchádzajúcemu videu. Náhodné poradie by tieto efekty eliminovalo, avšak pre 6 videí by vyžadovalo výrazne väčší počet respondentov.

\textbf{Jeden prezentujúci.} Všetky videá nakrútila jedna osoba, čo obmedzuje zovšeobecniteľnosť výsledkov --- vplyv dimenzií môže byť čiastočne špecifický pre konkrétneho prezentujúceho. Pre potvrdenie záverov by boli potrebné videá s viacerými prezentujúcimi.

\textbf{Málo videí pre validáciu modelu.} Spearmanova korelácia počítaná z $n = 6$ videí je citlivá na jednotlivé body. Bootstrap intervaly spoľahlivosti pri tejto veľkosti vzorky neposkytujú konfirmačnú informáciu. Validácia na nezávislom vzorke reálnych prezentácií je predmetom ďalšieho výskumu.
